\documentclass{article}
\usepackage{amsmath,amsthm}
\newtheorem{theorem}{Theorem}
\begin{document}
\title{When old evidence can certify clearance after refitting}
\maketitle

\paragraph{Scope.}
A civil warehouse controller considers admitting a vehicle to a zone. At a
decision cut it needs every occupied sensor sampling location to have
a certified clearance lower bound above the required margin. In the
development fixture the samples, associations and zone membership are fixed.
A calibration fit changes how two raw features jointly bound clearance.
A separately stipulated error envelope and aging debit turn that score
into a lower bound; regression does not verify either physical premise.
This model does not establish any sensor's accuracy or safe hardware
behavior.

\paragraph{A computed refit.}
The two models are fit by exact least squares to full-rank sets of four
authored labeled points. The old learned score is $2+x$; the updated score
is $x+y$. The two retained clouds are
$A=\{(1,2),(2,1)\}$ and $B=\{(1,1),(2,2)\}$.
They have identical feature marginals and identical old-model minima of
three. Both pass the threshold $2.5$ with stipulated error debit $0.25$.
The new minima are three and two, so only $A$ passes. If the first
capsule incurs additional age debit, the bound lasts through precisely
$0.25$ under the stated margin. If zone geometry or sample associations
changed, recomputing the old cloud's score is not that new decision.

\begin{theorem}
Fix a convex feature domain and any new lower-bound score $f$ on it.
The minimum score on every nonempty finite cloud is determined by its
convex hull, and equals the minimum score on the vertices, exactly
when $f$ is quasiconcave: throughout every line segment it is no lower
than the lesser endpoint score.
\end{theorem}
\begin{proof}
Quasiconcavity extends to finite convex combinations by induction.
The score on each interior sample is therefore at least the minimum
over the vertices. All vertices are samples; equality follows.
Conversely, take a point on a segment with score below both endpoints
and compare the two endpoints alone with the endpoints plus that point.
The hulls are identical; the minimum scores are not.
\end{proof}

Every affine fit satisfies the criterion. A non-affine fit requires a
test of its complete effective lower-bound score, including any
feature-dependent error debit. For $f(x,y)=x^2+y^2$ the sample $(0,0)$
changes the minimum of the cloud $\{(-1,0),(1,0)\}$ without changing
the hull. No hull-only re-evaluation can preserve both verdicts in
that extended model class. A convex hull summarizes a fixed collection
of paired feature vectors; it does not grant authorization, supply missing
features, re-associate samples, or preserve exact per-subzone minima
after splitting the admissible scope.

\paragraph{A received change is not necessarily used.}
The authored legacy consumer is given the current, correctly fitted model
and retains the old executed admission stamp; it decides using the old
score instead of the received update. With correct delivery and
interpretation held fixed, that mistake admits the $B$ case that the new
law leaves uncertified. This is an explicit omission control, not a
claim that any pinned fleet planner omits a valid closure. Reapplying the
new law over the retained hull ties the raw-data and matched incremental
checks. This tests certificate reuse under a declared refit family; there
is no novelty or architecture-disposition claim.
\end{document}
